\section{Introduction}
\label{sec:intro}

Enterprises accumulate process regulations, failure analyses, and operational
lessons as unstructured documents on shared drives. Answering a concrete
engineering question --- \emph{why did edge delamination appear in an
electroluminescence inspection of a laminated photovoltaic module?} ---
requires locating the one report among thousands that records the decision and
reading it. Retrieval-augmented generation (RAG) addresses the location step by
embedding documents and the query into a shared vector space
\cite{rag-lewis,dpr}, and agentic stacks let a large language model orchestrate
retrieval tools iteratively \cite{mcp,singh-agentic-rag}. The dominant ranking
signal, dense cosine similarity, measures \emph{topic relatedness} rather than
\emph{evidence presence} \cite{siren,gao-ragsurvey}.

Three actions this task requires are addressed by three research communities,
each solving its own and leaving the others' problems intact. Retrieval
research optimizes \emph{finding}: sparse, dense and hybrid rankers, and more
recently graph-based and hierarchical indexes, are judged by whether a relevant
passage appears in a top-$k$ list. Knowledge management addresses
\emph{organizing}: classification schemes, controlled vocabularies and
lessons-learned systems organize how documents are arranged and reused, but do
not close the loop to an answer a machine must justify. Document AI
addresses \emph{reading}: layout-aware parsing recovers text, tables and
formulas from scanned pages \cite{mineru}, but a parsed document still has to
be placed, indexed, governed and trusted. In deployment the three concerns
compound.

The object all three actions govern is \emph{content}: a knowledge-base
management system expresses that position as enforceable contracts rather than
as good practice. Three commitments follow, each implemented and measured here.
First, \emph{content decides placement}: library membership, tags, descriptions
and graph edges are derived from what a document says, not from the folder it
arrived in, and the resulting catalog is the first retrieval instrument.
Second, \emph{content decides admissibility}: no similarity score, however
high, promotes a candidate into an answer, because the platform must itself
fetch the candidate's content and judge it to contain direct evidence. Third,
\emph{absence of evidence is an answer}: when verification eliminates every
candidate, the system returns an explicit not-found report rather than the
closest text. Claims are therefore auditable, with every placement, verdict,
cut and lifecycle operation on record.

QDCVR (Query-Driven Content-Verified Retrieval, v2.3), the self-hosted
knowledge-base management system built here,
implements this position end to end on a deployed industrial corpus and on
acquisition and functional campaigns. Six functional subsystems carry the
platform: a quality-gated parsing and ingestion pipeline, knowledge
organization with content routing, tags and a document graph, life-cycle
management with three-source consistency audits, an operational-knowledge store
with credibility tiers, a multi-lane content-verified retrieval layer, and a
tool interface exposing the platform as 95 Model Context Protocol tools
\cite{mcp}; the system runs on a single workstation. On the acquisition side,
of 100 papers submitted to a 100-document campaign, 21 never became
retrievable, and the
design decision that matters is that a document which fails a gate is
reported rather than
stored as an empty document (Section~\ref{sec:e1}). On the retrieval side, the
content-verification gate is a better filter than a ranker: reusing the
\emph{same} verdicts as a re-ranking signal lowers Hit@5 from 0.94 to 0.44,
while using them as a binary filter restores 0.94 (sign test $p = 0.008$),
because same-domain score inflation makes non-gold candidates statistically
indistinguishable from gold ones (Mann--Whitney $p = 0.26$)
(Section~\ref{sec:eval-verify}).

\noindent Contributions:
\begin{enumerate}
\item \textbf{A content-governed knowledge-base management architecture}
(Sections~\ref{sec:system}--\ref{sec:implementation}): requirements derived
from deployment, six design principles enforced in code, a four-layer
architecture, and an explicit data model whose contracts --- parse gates,
placement rules, three-source consistency, promotion conditions, task
semantics --- are stated as obligations rather than conventions.
\item \textbf{A quality-gated acquisition pipeline with a measured failure
budget} (Section~\ref{sec:acquisition}, E1): multi-format parsing with OCR,
exact and near-duplicate detection, structure-aware splitting of oversized
documents with part addressing, content-derived descriptions and tags, and an
index post-condition; evaluated on a 100-document campaign (79 retrievable,
a 21\% retrievability gap whose cause the campaign record does not
attribute) and on a cross-lingual corpus.
\item \textbf{Content-routed organization and life-cycle management}
(Section~\ref{sec:organization}, E2): placement by content, catalog and
metadata registries with progressive disclosure, tag governance, duplicate
hygiene, a document graph with an explicitly restricted role, an experience
store with credibility tiers, and lifecycle operations certified by a
read-only audit (322 documents in five libraries).
\item \textbf{Multi-channel content-verified retrieval} (Section~\ref{sec:lanes},
E4--E5): three lanes over one tool surface under enforced budgets, a
verification gate served by a resident judging worker, and an honest
abstention path --- compared against nine retrieval configurations, two of
them derived offline from the same logged candidate pools, and against a
bare agent and a pure-dense controller. What the ablation measures is where
the verdict belongs, and what the platform delivers is an auditable record of
every candidate and every rejection, not a recall gain.
\item \textbf{A verifier-failure analysis and its role decomposition}
(Section~\ref{sec:eval-verify}, E4, E6): same-domain inflation quantified
inside a deployed pipeline, the filter/re-ranker asymmetry, a threshold that
never abstains out of corpus, and robustness of the failure to paraphrase, to
a second verification signal and to a second corpus.
\item \textbf{An operational account of the system} (Sections~\ref{sec:impl}
and~\ref{sec:discussion}, E7): resident-worker and subprocess engineering,
engine-agnostic harness integration with cost measured across three engines,
a catalogue of failure modes with their detection symptoms, and the
non-functional behaviour observed under concurrent writes.
\end{enumerate}

A demo describing the platform's organization and retrieval layers was
published as a CIKM~2026 demo \cite{qdcvr-demo}. This paper extends it with
four components the demo does not contain: the acquisition pipeline with its
measured retrievability gap and cross-lingual evaluation; life-cycle
management and integrity auditing as a designed subsystem rather than a
maintenance procedure; the operational-knowledge store with its credibility
tiers and distillation loop; and the verifier-failure analysis with the
filter/re-ranker decomposition. The demo also reports no negative result, which
is the centrepiece here.

\subsection{Research questions and reading guide}
\label{sec:rq}

The paper is organized as a design-to-evidence arc around five research
questions, answered by one evidence group each and reported in one place
(Table~\ref{tab:ataglance}).

\begin{description}
\item[RQ1 (Enforceable contracts).] How can the three content commitments be
turned into system contracts that hold across retrieval strategies and
independent agent harnesses, rather than into prompt conventions?
\item[RQ2 (Acquisition and organization at scale).] Do quality-gated
acquisition and content-routed organization produce collections that are
usable and internally consistent when they grow, and what does the gate
reject?
\item[RQ3 (Content-verified retrieval).] Does staged, multi-channel retrieval
with content verification outperform one-shot vector top-$k$ and unverified
fusion --- and by which mechanism?
\item[RQ4 (Cost and benefit of verification).] What does the gate cost and
buy, and is its failure mode structural rather than an artefact of one
checkpoint, domain or language?
\item[RQ5 (Operational contract).] Where do honest abstention, latency, cost
and stability actually live in a running system?
\end{description}

\begin{table*}[t]
\centering
\footnotesize
\setlength{\tabcolsep}{3pt}
\caption{This paper at a glance: each research question, the evidence group
that answers it, and the headline finding.}
\label{tab:ataglance}
\begin{tabular}{p{0.8cm}p{3.7cm}p{2.9cm}p{4.9cm}}
\toprule
RQ & Question & Evidence & Headline finding \\
\midrule
RQ1 & Are the content commitments enforceable contracts? &
Architecture, interface, implementation
(\S\ref{sec:system}--\S\ref{sec:implementation}) & Yes: gates, placement
rules, lane whitelists and fail-closed verdicts make them properties of the
enforcement layer \\
RQ2 & Do gated acquisition and content routing scale? &
E1, E2, E8 (\S\ref{sec:acquisition}, \S\ref{sec:organization}) & Partly: routing
and audits hold at 322 documents, while 21 of 100 submitted papers never
became retrievable and the record does not attribute the gap. Routing is
measured as coverage rather than accuracy, because every stored filename
carries the acquisition prefix (\S\ref{sec:e8}) \\
RQ3 & Does staged verified retrieval beat one-shot top-$k$? &
E4 (\S\ref{sec:eval-retrieval}) & Yes on reachability, through staged
recall plus content adjudication; two-stage retrieval is an accelerator, not
the contribution \\
RQ4 & What does the gate cost and buy, and is the failure structural? &
E4, E6 (\S\ref{sec:eval-verify}) & Costs $<$1\,s; as a re-ranker it loses
Hit@5 ($0.94\!\to\!0.44$), as a filter it is neutral ($0.94$,
$p=0.008$); the loss tracks confuser density, not checkpoint or domain \\
RQ5 & Where do honesty, latency and cost live? &
E5, E7 (\S\ref{sec:eval-agent}, \S\ref{sec:implementation}) & Honesty lives
in the reading and answering layer, never in the threshold; latency is
dominated by LLM turns, and engine choice moves cost by two orders of
magnitude \\
\bottomrule
\end{tabular}
\end{table*}

The remainder of the paper follows the arc.
Section~\ref{sec:related} positions the work against knowledge management,
document parsing, retrieval, verification and agent protocols.
Sections~\ref{sec:system}--\ref{sec:organization} present the architecture,
the acquisition pipeline and the organization subsystem, answering RQ1 and
RQ2. Section~\ref{sec:lanes} presents the retrieval lanes and the verification
gate, walking one real question end to end, and states what the retrieval layer
claims (RQ3). Section~\ref{sec:implementation} reports the engineering that
makes the contracts enforceable. Section~\ref{sec:eval} carries the eight
evidence groups (RQ4--RQ5).
Sections~\ref{sec:discussion}--\ref{sec:conclusion} distil the lessons and
state the threats to validity. Table~\ref{tab:modules} maps system functions
to components and evidence; Table~\ref{tab:ataglance} maps research questions
to evidence.
